% --- Archon physics formalization source begin ---
% archon:physics
% archon:covers PhyXMiniProblems/problem_phyx_mini_0083.lean
% archon:source-report reports/phyx_mini/problem_phyx_mini_0083.source.json

\chapter{Physics problem phyx\_mini\_0083}
\label{ch:PhyXMiniProblems_problem_phyx_mini_0083}

\paragraph{Problem source.}
\#\# Physical scenario

Figure gives \textbackslash{}( \textbackslash{}alpha \textbackslash{}) versus the sine of the angle \textbackslash{}( \textbackslash{}theta \textbackslash{}) in a single-slit diffraction experiment using light of wavelength 610\textbackslash{},\textbackslash{}text\{nm\}. The vertical axis scale is set by \textbackslash{}( \textbackslash{}alpha\_s = 12\textbackslash{},\textbackslash{}text\{rad\} \textbackslash{}).

\#\# Question

What is the slit width?

\#\# Answer choices

- A: A: 2.10μm

- B: B: 2.50μm

- C: C: 2.33μm

- D: D: 1.95μm

\#\# Figure caption (auxiliary; use the image as primary evidence)

The image is a graph with the following elements:

- **Axes:**
  - The x-axis is labeled "sin θ" with values ranging from 0 to 1.
  - The y-axis is labeled "α (rad)" with a specific point labeled as "αₛ".

- **Line:**
  - A bold diagonal line runs from the origin (0, 0) to the point (1, αₛ) at the top right corner of the grid.

- **Grid:**
  - The background has a regular grid with vertical and horizontal lines, creating smaller squares.

The graph depicts a linear relationship between "sin θ" and "α" in radians, illustrating that as "sin θ" increases from 0 to 1, "α" increases linearly to "αₛ".

\#\# Dataset metadata

Wave Optics; Physical Model Grounding Reasoning, Multi-Formula Reasoning

\paragraph{Recorded answer/context.}
C: C: 2.33μm

\paragraph{Figure/image path.}
/root/proposal\_for\_physic/science-mango/phyx\_mini\_run/phyx\_data/test\_image/83.png

\paragraph{Formalization target.}
create a compiling Lean file with sorry bodies at `PhyXMiniProblems/problem\_phyx\_mini\_0083.lean`. The Lean declarations must preserve the physical quantities, dimensions or dimensional roles, figure labels, governing-law hypotheses, and final relation expressed by this problem.
Use Mathlib/Physlib names found through LeanExplore where available. If a physics API is missing, introduce faithful local abstractions rather than scalar placeholder aliases.

\begin{theorem}[Physics formalization target]
\label{thm:physics:phyx_mini_0083:target}
\lean{PhyXMiniProblems.ProblemPhyXMini0083.slitWidth_rounds_to_recordedAnswerC}
\uses{def:physics:phyx-mini-0083:phyxminiproblems-problemphyxmini0083-micrometersvalue, def:physics:phyx-mini-0083:phyxminiproblems-problemphyxmini0083-singleslitexperiment, def:physics:phyx-mini-0083:phyxminiproblems-problemphyxmini0083-alphaversussinethetagraph, def:physics:phyx-mini-0083:phyxminiproblems-problemphyxmini0083-satisfiessingleslitphaselaw, def:physics:phyx-mini-0083:phyxminiproblems-problemphyxmini0083-hasstatedwavelengthandgraphreadouts, def:physics:phyx-mini-0083:phyxminiproblems-problemphyxmini0083-answerslitwidthhundredthsmicrometers}
For wavelength \(610\,\mathrm{nm}\) and the linear graph
\(\alpha=12\sin\theta\), the single-slit phase law determines the slit width;
its micrometre readout rounds to \(2.33\,\mu\mathrm m\), answer C.
\end{theorem}
\begin{proof}
The graph slope is \(12\) radians per unit of \(\sin\theta\).  The phase law
\(\alpha=\pi a\sin\theta/\lambda\) therefore gives
\(a=12\lambda/\pi\).  Insert \(\lambda=610\,\mathrm{nm}\) and convert
nanometres to micrometres, obtaining \(a=7.32/\pi\,\mu\mathrm m\).
Certified rational bounds for \(\pi\) place
\(100a\) in the half-open rounding interval for the integer \(233\).
Consequently the nearest-integer readout equals the hundredths-of-a-
micrometre value attached to answer C.
\end{proof}

% --- Archon declaration topology begin ---
\section{Declaration topology}
\begin{definition}[Dim Length]
  \label{def:physics:phyx-mini-0083:phyxminiproblems-problemphyxmini0083-dimlength}
  \lean{PhyXMiniProblems.ProblemPhyXMini0083.DimLength}
  A signed physical length represented independently of a chosen unit system
\end{definition}
\begin{proof}
  This is the declared model definition of Dim Length.
\end{proof}
\begin{definition}[length Value In]
  \label{def:physics:phyx-mini-0083:phyxminiproblems-problemphyxmini0083-lengthvaluein}
  \lean{PhyXMiniProblems.ProblemPhyXMini0083.lengthValueIn}
  \uses{def:physics:phyx-mini-0083:phyxminiproblems-problemphyxmini0083-dimlength}
  Read a physical length as a real scalar in the specified length unit
\end{definition}
\begin{proof}
  This is the declared model definition of length Value In.
\end{proof}
\begin{definition}[meters Value]
  \label{def:physics:phyx-mini-0083:phyxminiproblems-problemphyxmini0083-metersvalue}
  \lean{PhyXMiniProblems.ProblemPhyXMini0083.metersValue}
  \uses{def:physics:phyx-mini-0083:phyxminiproblems-problemphyxmini0083-dimlength, def:physics:phyx-mini-0083:phyxminiproblems-problemphyxmini0083-lengthvaluein}
  The metre readout used in the dimensionless single-slit phase law
\end{definition}
\begin{proof}
  This is the declared model definition of meters Value.
\end{proof}
\begin{definition}[nanometers Value]
  \label{def:physics:phyx-mini-0083:phyxminiproblems-problemphyxmini0083-nanometersvalue}
  \lean{PhyXMiniProblems.ProblemPhyXMini0083.nanometersValue}
  \uses{def:physics:phyx-mini-0083:phyxminiproblems-problemphyxmini0083-dimlength, def:physics:phyx-mini-0083:phyxminiproblems-problemphyxmini0083-lengthvaluein}
  The nanometre readout used for the stated wavelength
\end{definition}
\begin{proof}
  This is the declared model definition of nanometers Value.
\end{proof}
\begin{definition}[micrometers Value]
  \label{def:physics:phyx-mini-0083:phyxminiproblems-problemphyxmini0083-micrometersvalue}
  \lean{PhyXMiniProblems.ProblemPhyXMini0083.micrometersValue}
  \uses{def:physics:phyx-mini-0083:phyxminiproblems-problemphyxmini0083-dimlength, def:physics:phyx-mini-0083:phyxminiproblems-problemphyxmini0083-lengthvaluein}
  The micrometre readout used by the four answer choices
\end{definition}
\begin{proof}
  This is the declared model definition of micrometers Value.
\end{proof}
\begin{definition}[Single Slit Experiment]
  \label{def:physics:phyx-mini-0083:phyxminiproblems-problemphyxmini0083-singleslitexperiment}
  \lean{PhyXMiniProblems.ProblemPhyXMini0083.SingleSlitExperiment}
  \uses{def:physics:phyx-mini-0083:phyxminiproblems-problemphyxmini0083-dimlength}
  The monochromatic illumination and single slit. The wavelength is the stated 610 nm; the slit width is the physical quantity requested by the problem
\end{definition}
\begin{proof}
  This is the declared model definition of Single Slit Experiment.
\end{proof}
\begin{definition}[Alpha Versus Sine Theta Graph]
  \label{def:physics:phyx-mini-0083:phyxminiproblems-problemphyxmini0083-alphaversussinethetagraph}
  \lean{PhyXMiniProblems.ProblemPhyXMini0083.AlphaVersusSineThetaGraph}
  The plotted graph. Its horizontal coordinate is a scalar value of sin θ, and its vertical coordinate is the phase parameter α read in radians. The field alphaScaleRadians is the figure label αₛ at the top of the vertical axis
\end{definition}
\begin{proof}
  This is the declared model definition of Alpha Versus Sine Theta Graph.
\end{proof}
\begin{definition}[Satisfies Single Slit Phase Law]
  \label{def:physics:phyx-mini-0083:phyxminiproblems-problemphyxmini0083-satisfiessingleslitphaselaw}
  \lean{PhyXMiniProblems.ProblemPhyXMini0083.SatisfiesSingleSlitPhaseLaw}
  \uses{def:physics:phyx-mini-0083:phyxminiproblems-problemphyxmini0083-metersvalue, def:physics:phyx-mini-0083:phyxminiproblems-problemphyxmini0083-singleslitexperiment, def:physics:phyx-mini-0083:phyxminiproblems-problemphyxmini0083-alphaversussinethetagraph}
  The Fraunhofer single-slit phase-parameter law α = (π a / λ) sin θ on the displayed interval 0 ≤ sin θ ≤ 1. Metre readouts are used for both lengths, so their quotient is dimensionless. This governing law contains no numerical value for the requested slit width
\end{definition}
\begin{proof}
  This is the declared model definition of Satisfies Single Slit Phase Law.
\end{proof}
\begin{definition}[Has Stated Wavelength And Graph Readouts]
  \label{def:physics:phyx-mini-0083:phyxminiproblems-problemphyxmini0083-hasstatedwavelengthandgraphreadouts}
  \lean{PhyXMiniProblems.ProblemPhyXMini0083.HasStatedWavelengthAndGraphReadouts}
  \uses{def:physics:phyx-mini-0083:phyxminiproblems-problemphyxmini0083-nanometersvalue, def:physics:phyx-mini-0083:phyxminiproblems-problemphyxmini0083-singleslitexperiment, def:physics:phyx-mini-0083:phyxminiproblems-problemphyxmini0083-alphaversussinethetagraph}
  The problem-text and primary-figure readouts: λ = 610 nm, αₛ = 12 rad, and the bold line from (0, 0) to (1, αₛ). The universal equation records the displayed straight line throughout the horizontal range. No slit-width readout occurs in this predicate
\end{definition}
\begin{proof}
  This is the declared model definition of Has Stated Wavelength And Graph Readouts.
\end{proof}
\begin{definition}[Answer Choice]
  \label{def:physics:phyx-mini-0083:phyxminiproblems-problemphyxmini0083-answerchoice}
  \lean{PhyXMiniProblems.ProblemPhyXMini0083.AnswerChoice}
  Labels of the four multiple-choice answers printed with the problem
\end{definition}
\begin{proof}
  This is the declared model definition of Answer Choice.
\end{proof}
\begin{definition}[answer Slit Width Hundredths Micrometers]
  \label{def:physics:phyx-mini-0083:phyxminiproblems-problemphyxmini0083-answerslitwidthhundredthsmicrometers}
  \lean{PhyXMiniProblems.ProblemPhyXMini0083.answerSlitWidthHundredthsMicrometers}
  \uses{def:physics:phyx-mini-0083:phyxminiproblems-problemphyxmini0083-answerchoice}
  Each displayed slit-width choice, in hundredths of a micrometre. Thus the recorded value for choice C is 233 / 100 μm = 2.33 μm
\end{definition}
\begin{proof}
  This is the declared model definition of answer Slit Width Hundredths Micrometers.
\end{proof}
% --- Archon declaration topology end ---
% --- Archon physics formalization source end ---
